% R1 – System Architecture Design: network protocols used in the project
% Compile: pdflatex R1_protocols.tex   (or upload to Overleaf)
\documentclass[11pt,a4paper]{article}

\usepackage[margin=2.2cm]{geometry}
\usepackage[T1]{fontenc}
\usepackage{lmodern}
\usepackage{booktabs}
\usepackage{tabularx}
\usepackage{array}
\usepackage{enumitem}
\usepackage{xcolor}
\usepackage{listings}
\usepackage{hyperref}

\definecolor{codebg}{HTML}{F4F6F8}
\lstset{
  basicstyle=\ttfamily\small,
  backgroundcolor=\color{codebg},
  frame=single, rulecolor=\color{gray!40},
  columns=fullflexible, keepspaces=true,
  breaklines=true, xleftmargin=4pt, xrightmargin=4pt
}
\newcolumntype{L}{>{\raggedright\arraybackslash}X}
\setlist{nosep, leftmargin=*}
\newcommand{\code}[1]{\texttt{#1}}

\title{\textbf{R1 -- System Architecture Design}\\[2pt]
\large Computer Network Protocols Used in the Simple Chat Application}
\author{{[Team name]} \\ {[Member names]}}
\date{}

\begin{document}
\maketitle

\section{Overview: the protocol stack}

The system is \textbf{client--server}: two or more browser clients on separate
computers connect to one Node.js server, and the server decides who receives
each message. Every chat message travels through the following stack:

\begin{lstlisting}
Socket.IO events      <- our application protocol (register, private_msg, group_msg, ...)
WebSocket (RFC 6455)  <- full-duplex frames on one long-lived connection
HTTP/1.1              <- used once per connection, only to upgrade to WebSocket
TCP                   <- reliable, ordered byte stream to the server's port 3000
IPv4                  <- addresses each computer on the LAN / Wi-Fi
Wi-Fi / Ethernet
\end{lstlisting}

\section{Core protocols}

\subsection{IPv4 (Network layer)}
\begin{description}
  \item[What:] The addressing protocol that identifies each computer on the network.
  \item[Why:] Our clients run on different physical computers, so they need to
    find the server machine. IPv4 is used by every LAN, Wi-Fi network and phone hotspot.
  \item[How:]
    \begin{itemize}
      \item The server listens on \code{0.0.0.0} (all network interfaces), not
        \code{localhost}, so other machines can reach it.
      \item At startup the server prints its IPv4 addresses.
      \item Users type \code{http://<server IP>:3000} on the login screen. Nothing is hard-coded.
    \end{itemize}
\end{description}

\subsection{TCP (Transport layer)}
\begin{description}
  \item[What:] A connection-oriented transport protocol. It uses a 3-way
    handshake, delivers data reliably and in order, and retransmits lost data.
  \item[Why:] In a chat, a lost or reordered message is a visible error. TCP
    guarantees both delivery and order; UDP guarantees neither. Chat traffic is
    small, so TCP's overhead doesn't matter.
  \item[How:]
    \begin{itemize}
      \item Each client keeps \textbf{one TCP connection} to port 3000 for its whole session.
      \item The connection's lifetime is the user's online status. When a tab
        closes, TCP closes and the server receives \code{disconnect}. It then
        frees the user's name, removes the user from all groups and broadcasts
        the updated lists.
    \end{itemize}
\end{description}

\subsection{HTTP/1.1 (Application layer, used once per connection)}
\begin{description}
  \item[What:] The web's request/response protocol.
  \item[Why:] A WebSocket connection is opened with a standard HTTP request, so
    it passes through firewalls like normal web traffic.
  \item[How:]
    \begin{itemize}
      \item The browser sends one \code{GET /socket.io/?EIO=4\&transport=websocket}
        with the header \code{Upgrade: websocket}.
      \item The server answers \code{101 Switching Protocols}, and the same TCP
        connection becomes a WebSocket.
      \item \textbf{Every other HTTP path returns 404.} There is no REST API for
        sending or fetching messages. This shows that chat uses sockets only (R2).
    \end{itemize}
\end{description}

\subsection{WebSocket (RFC 6455)}
\begin{description}
  \item[What:] A full-duplex protocol where both sides can send frames at any
    time over one persistent connection.
  \item[Why:] With HTTP polling, the client must repeatedly ask ``anything new?''.
    That adds delay and wastes requests. With WebSocket, the server
    \textbf{pushes} a message the moment it arrives.
  \item[How:]
    \begin{itemize}
      \item The client forces \code{transports: ["websocket"]}, so no polling is used.
      \item Every event is one frame, visible in DevTools $\rightarrow$ Network
        $\rightarrow$ WS, for example:
\begin{lstlisting}
42["private_msg",{"to":"C2","text":"hi"}]
\end{lstlisting}
        Here \code{4} is the Engine.IO message type and \code{2} is the Socket.IO event type.
      \item File data is sent as \textbf{binary frames}.
    \end{itemize}
\end{description}

\subsection{Socket.IO (our application protocol)}
\begin{description}
  \item[What:] A protocol on top of WebSocket that adds named events, rooms,
    acknowledgements and a heartbeat. The set of events and their payloads is
    our own message format, documented in \code{docs/PROTOCOL.md}.
  \item[Why:] Browsers cannot open raw TCP sockets. Socket.IO works in the
    browser and still gives one long-lived socket per client. Its features match
    our requirements directly:
    \begin{itemize}
      \item \textbf{Named events} give message framing: \code{register},
        \code{private\_msg}, \code{create\_group}, \code{join\_group}, \code{group\_msg}.
      \item \textbf{Rooms} implement groups: \code{io.to("group:g1")} reaches
        only sockets in that room (R11).
      \item \textbf{Acknowledgements} give request/response over the socket (games, file chunks).
      \item The \textbf{heartbeat} (a ping every 25\,s, with a 20\,s timeout)
        detects connections that dropped without closing.
    \end{itemize}
  \item[How:]
    \begin{itemize}
      \item \textbf{Private message (R7):} the server emits only to the
        receiver's socket, plus an echo to the sender.
      \item \textbf{Group message (R11):} the server first checks membership,
        then emits to the group's room only.
      \item The \textbf{sender's identity comes from the name bound to the socket
        at \code{register}}, never from the message, so nobody can impersonate another user.
      \item The server logs every received event
        (\code{[recv] C1 -> private\_msg \{...\}}) as live evidence that all chat
        goes through the socket.
    \end{itemize}
\end{description}

\section{Protocols used by the special features}

{\small
\renewcommand{\arraystretch}{1.35}
\begin{tabularx}{\textwidth}{@{}>{\raggedright\arraybackslash}p{2.3cm} L L >{\raggedright\arraybackslash}p{6.4cm}@{}}
\toprule
\textbf{Feature} & \textbf{Protocols} & \textbf{Why} & \textbf{How} \\
\midrule
\textbf{AI word-guessing game}
& Socket.IO (the AI is another client) + \textbf{HTTPS} to the AI APIs
& Every AI message must go through the socket like a human's. HTTPS encrypts
  the requests and keeps the API keys secret.
& The AI bot is a separate program that logs in as user \code{AI} using a
  secret token. The server forwards \code{/guess} commands to it as
  \code{bot\_request} (with an ack). The bot joins the group itself and sends
  hints as normal \code{group\_msg}. It calls Claude / ChatGPT / Gemini over
  HTTPS. The answer and the keys stay only on the bot. \\
\midrule
\textbf{Rock-paper-scissors game}
& Our own \textbf{application-level protocol} on Socket.IO: request/ack +
  a server-side state machine
& Both players must choose ``at the same time'', and neither may see the
  other's move first.
& Every command (\code{game\_challenge}, \code{game\_respond}, \code{game\_move},
  \code{game\_quit}) gets an ack \code{\{ok, code\}}. The server tracks
  \code{pending} $\rightarrow$ \code{playing} $\rightarrow$ \code{finished}. It
  keeps each move hidden and tells the opponent only that a move was made
  (\code{moved}). Once both have chosen, it reveals both moves in one
  \code{round} update. Updates go only to the 2 players. A disconnect ends the
  game. Nobody can cheat by reading WebSocket frames. \\
\midrule
\textbf{File / picture sending}
& Binary Socket.IO frames + \textbf{chunking} + \textbf{stop-and-wait ACK}
& Sending a whole file in one event would exceed Socket.IO's 1\,MB message
  limit and show no progress.
& Files ($\le$ 5\,MB) are split into 64\,KB chunks with framing
  \code{\{id, seq, data\}}. The sender waits for each chunk's ack before
  sending the next. This is flow control at the application level, like TCP
  ACKs. The server checks who owns the transfer, the chunk order and the size,
  then relays each chunk only to the receiver or group members. It never
  stores the file. \\
\midrule
\textbf{Discord bridge}
& \textbf{Discord Gateway} (WebSocket over TLS) + \textbf{Discord REST API} (HTTPS)
& The Gateway delivers Discord messages in real time; REST is how Discord
  accepts new posts.
& Our server acts as a \emph{client} of Discord. It receives
  \code{MessageCreate} events through the Gateway and forwards them to the
  \code{\#discord} group room. When a member posts in that group, it sends the
  message to Discord with \code{POST /channels/\{id\}/messages}. Inside our app,
  everything still travels over Socket.IO. \\
\bottomrule
\end{tabularx}
}

\section{Summary}

Our chat uses \textbf{Socket.IO events over WebSocket over TCP over IPv4}. HTTP
is used only once per connection to upgrade to WebSocket. \textbf{TCP}
guarantees messages arrive complete and in order, \textbf{WebSocket} lets the
server push messages instantly, and \textbf{Socket.IO} rooms and events let the
server alone decide who receives each message. The special features reuse the
same socket. They add HTTPS for AI calls, an application-level game protocol,
chunked binary transfer with ACKs, and Discord's Gateway and REST APIs.

\end{document}
